\documentclass[10pt,letterpaper]{article}
\usepackage[margin=1in]{geometry}
\usepackage{amsmath,amssymb}
\usepackage{booktabs}
\usepackage{graphicx}
\usepackage{hyperref}
\usepackage{xcolor}
\usepackage{multirow}
\usepackage{array}

\title{Why Temporal Abstraction Fails in DreamerV3: \\A Single-Environment Diagnostic Analysis of Variable-Rate Imagination}

\author{Anonymous Authors}
\date{}

\begin{document}
\maketitle

\begin{abstract}
Temporal abstraction promises more efficient planning in world-model reinforcement learning, yet variable-rate imagination consistently degrades performance in DreamerV3. We present a rigorous factorial diagnosis within Walker Walk that identifies the specific causal chain: a trained duration-conditioning embedding (dt-embedding) amplifies RSSM fidelity gaps through observation-decoder gradients, producing \emph{conditioning-induced value overestimation} (CIVO)---an accelerating positive feedback loop where the critic systematically overestimates returns for temporally-abstracted transitions. Our normalized value-to-evaluation ratio (NVE) diagnostic achieves perfect classification (5/5 accuracy, AUROC 0.98) in Walker Walk, though it does not generalize to Cheetah Run or Cartpole Swingup where degradation proceeds through distinct mechanisms. Gradient-level intervention (Full Selective Gradient Stopping) reduces NVE to 1.00$\times$, confirming the failure is observation-decoder-driven rather than reward-driven. We contrast our finding with TAWM (Nhu et al., ICML 2025), where dt-conditioning succeeds in model-predictive control but fails in imagination-based actor-critic training, providing the first mechanistic explanation for this discrepancy. Our results demonstrate that temporal abstraction fails universally across three environments but through environment-specific mechanisms, complicating one-size-fits-all solutions.
\end{abstract}

%=============================================
\section{Introduction}
%=============================================

Temporal abstraction---the ability to reason and plan over variable time scales---is a foundational goal in reinforcement learning \cite{sutton1999options,vezhnevets2017feudal}. World-model RL architectures such as DreamerV3 \cite{hafner2023dreamerv3} offer a natural substrate for temporal abstraction: one can modify the imagination step duration (dt) without altering environment interaction. Yet practical attempts at variable-rate imagination consistently degrade performance. Why?

We present a systematic factorial diagnosis within DreamerV3's imagination-based actor-critic framework, focusing on Walker Walk as our primary environment. Our investigation reveals a specific causal chain:
\begin{enumerate}
    \item A \emph{trained} dt-embedding enables the critic to exploit precise but erroneous temporal predictions from the RSSM at dt$>$1.
    \item The RSSM's fidelity gap at longer temporal intervals creates systematic value overestimation, amplified through observation-decoder gradients (93\% of dt-embedding gradient magnitude).
    \item This conditioning-induced value overestimation (CIVO) forms a positive feedback loop: overestimated values $\rightarrow$ exploitative actor $\rightarrow$ degraded evaluation $\rightarrow$ further relative overestimation.
\end{enumerate}

We further show that this mechanism is environment-specific: Cheetah Run and Cartpole Swingup both degrade with dt$>$1 transitions but through qualitatively different mechanisms (learning signal reduction and value suppression, respectively), without the accelerating CIVO feedback loop.

Our contributions:
\begin{itemize}
    \item The first rigorous factorial diagnosis revealing the specific causal chain of temporal abstraction failure in DreamerV3's imagination-based actor-critic training.
    \item The NVE diagnostic methodology---a normalized overestimation metric achieving perfect classification in Walker Walk (proposed for broader adoption with documented limitations).
    \item Gradient-level causal evidence: Full Selective Gradient Stopping (Full-SGS) on observation decoders reduces NVE to 1.00$\times$, proving failure is observation-decoder-driven (not reward-driven, since reward gradient $= 0$ by architecture).
    \item A mechanistic explanation for why dt-conditioning succeeds in MPC (TAWM, ICML 2025) but fails in imagination-based training.
\end{itemize}

%=============================================
\section{Background and Related Work}
%=============================================

\paragraph{Temporal abstraction in RL.}
The options framework \cite{sutton1999options} and its descendants \cite{bacon2017optioncritic,vezhnevets2017feudal} enable agents to operate at multiple time scales. In model-free settings, action repetition \cite{lakshminarayanan2017dynamic} and FiGAR \cite{sharma2017learning} provide simpler temporal abstraction mechanisms.

\paragraph{World-model RL.}
DreamerV3 \cite{hafner2023dreamerv3} learns a Recurrent State-Space Model (RSSM) and trains actor-critic networks entirely within imagined trajectories. The RSSM's latent dynamics model provides a natural interface for variable-rate prediction.

\paragraph{Variable-rate planning in world models.}
TAWM (Nhu et al., ICML 2025) demonstrates successful dt-conditioning in TD-MPC2's model-predictive control setting. Hierarchical approaches (THICK, Hieros) avoid flat dt-conditioning by structuring temporal abstraction through explicit hierarchy. Our work identifies why the simpler flat conditioning approach fails specifically in imagination-based actor-critic training.

\paragraph{Value overestimation in MBRL.}
Model exploitation and value overestimation are known challenges in model-based RL \cite{janner2019mbpo}. Lu et al.\ (2018) formalize delusional bias from policy-class constraints---conceptually distinct from our CIVO, which arises from model imprecision amplified by learned conditioning.

%=============================================
\section{Method}
%=============================================

\subsection{Architecture}

We use DreamerV3 as the base architecture with its RSSM world model. Variable-rate imagination is implemented via a \emph{dt-embedding}---a learned embedding vector concatenated to the RSSM hidden state that conditions predictions on the intended temporal skip duration.

\subsection{Factorial Experiment Design}

We design a controlled factorial experiment isolating three factors:
\begin{itemize}
    \item \textbf{dt-embedding status}: trained / untrained (random fixed) / absent (dim=0) / frozen
    \item \textbf{Temporal subsampling}: ON (dt$>$1 transitions in replay) / OFF (dt=1 only)
    \item \textbf{Imagination dt}: dt=1 / dt=2 / dt=4
\end{itemize}

The critical controlled comparison is between configurations that differ only in dt-embedding status while holding all other factors constant.

\subsection{Diagnostic Metric: Normalized Value-to-Evaluation Ratio (NVE)}

We define NVE as:
\begin{equation}
    \text{NVE}(t) = \frac{V_{\text{exp}}(t) / E_{\text{exp}}(t)}{V_{\text{baseline}}(t) / E_{\text{baseline}}(t)}
\end{equation}
where $V$ denotes mean predicted value and $E$ denotes evaluation return. NVE $> 1$ indicates relative overestimation; NVE $> 1.5\times$ at training steps $\geq 35$K (where baseline v/e has converged above 0.24) is our proposed diagnostic threshold for pathological overestimation. The 35K lower bound avoids false positives from small-denominator instability in early training.

\subsection{Environments}

We evaluate on three DeepMind Control Suite environments:
\begin{itemize}
    \item \textbf{Walker Walk}: 24D state, complex locomotion reward (primary environment).
    \item \textbf{Cheetah Run}: Velocity-based dense reward, smoother state dynamics.
    \item \textbf{Cartpole Swingup}: 4D state, sparse initial reward structure.
\end{itemize}

%=============================================
\section{Experiments}
%=============================================

\subsection{Walker Walk: CIVO Mechanism Identification}

\subsubsection{Trained dt-Embedding Reliably Triggers Pathological Overestimation}

Table~\ref{tab:008d_crossseed} presents multi-seed results for the CIVO-triggering configuration (trained dt-embedding, subsample ON, dt=2 imagination).

\begin{table}[h]
\centering
\caption{Walker Walk 008d configuration (trained dt-emb, subsample ON, dt=2 imag): cross-seed comparison.}
\label{tab:008d_crossseed}
\begin{tabular}{lccc}
\toprule
Metric & Seed 0 (008d) & Seed 1 (009e) & Seed 2 (009b) \\
\midrule
Peak NVE (35K+) & 2.09$\times$@45K & 4.12$\times$@55K & 1.62$\times$@45K \\
Collapse? & Yes@55K & Yes@55K & No (self-correct) \\
Final eval@105K & 82.7 (55K stop) & 131.1 & 267.9 \\
NVE trajectory & accelerating & accelerating & transient peak $\rightarrow$ decline \\
\bottomrule
\end{tabular}
\end{table}

All three seeds exceed NVE $> 1.5\times$ (100\% sensitivity), confirming trained dt-embedding reliably triggers pathological overestimation. However, collapse is stochastic: 2/3 seeds show catastrophic collapse while seed 2 self-corrects (NVE 2.05$\times$@25K $\rightarrow$ 0.38$\times$@105K). The bifurcation between sustained and transient overestimation depends on early training stochasticity.

\subsubsection{Causal Control: Removing dt-Embedding}

\begin{table}[h]
\centering
\caption{Walker Walk 008f configuration (dt\_embed\_dim=0, subsample ON): cross-seed comparison.}
\label{tab:008f_crossseed}
\begin{tabular}{lcc}
\toprule
Metric & Seed 0 (008f) & Seed 2 (009c) \\
\midrule
Peak NVE & 1.02$\times$@25K & 1.26$\times$@25K \\
Final eval@105K & 503 & 148 \\
Recovery pattern & Zigzag $\rightarrow$ 503 & Peak 359@75K $\rightarrow$ decline \\
vs baseline & $-$45\% & $-$84\% \\
\bottomrule
\end{tabular}
\end{table}

Without dt-embedding, NVE never exceeds 1.5$\times$ (0/2 seeds), confirming the causal role of trained dt-embedding in triggering CIVO. However, removing dt-embedding does not guarantee learning success---seed 2 (009c) shows severe performance degradation through a different mechanism (reward prediction quality collapses to 75$\times$ baseline MAE).

\subsubsection{Frozen dt-Embedding: Isolating Training as the Causal Factor}

Freezing the dt-embedding at random initialization (preventing gradient updates) produces a distinct failure mode: gradient starvation. The uninformative frozen embedding provides no useful conditioning to the RSSM, leading to systematic value underestimation (NVE 0.39--0.47$\times$), actor loss approaching zero (0.000179@105K), and unstable evaluation (80--238 range). This confirms that dt-embedding \emph{training}---not its mere presence---drives CIVO. The frozen condition serves as a factorial ablation point establishing: absent $\neq$ frozen $\neq$ trained.

Final evaluation: 204.2@105K (gate\_dt\_mean fixed at 2.0).

\subsubsection{Untrained dt-Embedding: Value Growth Suppression}

\begin{table}[h]
\centering
\caption{Walker Walk 008a (untrained dt-emb, subsample OFF, H=15): cross-seed summary.}
\label{tab:008a_crossseed}
\begin{tabular}{lcccc}
\toprule
Metric & Seed 0 & Seed 1 & Seed 2 & Mean \\
\midrule
Eval@105K & 778 & 700 & 620 & 699 \\
vs baseline & $-$15\% & $-$24\% & $-$32\% & $-$24\% \\
NVE@105K & 0.59$\times$ & 0.65$\times$ & 0.72$\times$ & 0.65$\times$ \\
Collapse? & No & No & No & 0/3 \\
\bottomrule
\end{tabular}
\end{table}

Untrained (random fixed) dt-embedding prevents CIVO across all seeds (3/3 NVE $< 1.0\times$, zero collapses). Mean performance reaches 76\% of baseline (699 vs 916), demonstrating structural value growth suppression without catastrophic failure.

\subsubsection{dt=1 Imagination: Partial Mitigation (C6 Partially Falsified)}

\begin{table}[h]
\centering
\caption{Walker Walk 008e\_v3 (trained dt-emb, dt=1 imagination, subsample ON): all seeds.}
\label{tab:008ev3_seeds}
\begin{tabular}{lcccc}
\toprule
Seed & Eval@105K & NVE@105K & Collapse? & Mechanism \\
\midrule
Seed 1 & 869.9 & 0.80$\times$ & No & --- \\
Seed 2 & 137.0 & 3.78$\times$ & Yes & Eval-driven (Pathway B) \\
Seed 3 & 830 & 0.78$\times$ & No & --- \\
Seed 4 & 790 & 0.82$\times$ & No & --- \\
\bottomrule
\end{tabular}
\end{table}

Restricting imagination to dt=1 while retaining trained dt-embedding yields a 25\% collapse rate (1/4 seeds). The collapsed seed (s2) exhibits dynamics qualitatively distinct from classic CIVO:

\begin{table}[h]
\centering
\caption{Pathway A (CIVO) vs Pathway B (non-CIVO) collapse signatures.}
\label{tab:pathways}
\begin{tabular}{lcc}
\toprule
Feature & 008d s0 (Pathway A: CIVO) & 008e\_v3 s2 (Pathway B) \\
\midrule
Value change & +130\% explosion & +13\% stable \\
Eval change & Post-decline $-$43\% & Pre-decline $-$81\% \\
NVE driver & Value (numerator) & Eval (denominator) \\
Actor entropy & $\downarrow$ ($-$3.47$\rightarrow$$-$4.88) & $\uparrow$ ($-$5.21$\rightarrow$$-$3.83) \\
Mechanism & Critic overestimate $\rightarrow$ exploit & Eval crash $\rightarrow$ value lag \\
\bottomrule
\end{tabular}
\end{table}

With 4 seeds, Pathway B collapse rate (25\%) is closer to seed variance than a systematic second mechanism. Honest reporting requires acknowledging $n=4$ provides insufficient statistical power to distinguish from baseline stochasticity.

\subsection{NVE Diagnostic Threshold Analysis}

\begin{table}[h]
\centering
\caption{NVE 1.5$\times$ threshold performance across environments (35K+ constraint).}
\label{tab:nve_crossenv}
\begin{tabular}{lccccccc}
\toprule
Environment & Trained & TP & FN & Untrained & TN & FP & Accuracy \\
\midrule
Walker Walk & 3 & 3 & 0 & 2 & 2 & 0 & 5/5 = 1.00 \\
Cheetah Run & 3 & 0 & 3 & 3 & 3 & 0 & 3/6 = 0.50 \\
Cartpole Swingup & 3 & 0 & 3 & 3 & 3 & 0 & 3/6 = 0.50 \\
\midrule
Cross-env total & 9 & 3 & 6 & 8 & 8 & 0 & 11/17 = 0.65 \\
\bottomrule
\end{tabular}
\end{table}

Within Walker Walk, the NVE $> 1.5\times$ threshold at 35K+ achieves perfect classification (AUROC $\approx 0.98$). Cross-environment, accuracy degrades to 0.65 (AUROC 0.72) due to complete failure in Cheetah Run and Cartpole Swingup, where degradation manifests without exceeding the threshold. The NVE diagnostic is validated as a Walker Walk--specific tool; broader adoption requires environment-specific calibration.

\subsection{Cross-Environment Generalization: Negative Results}

\subsubsection{Cheetah Run: Degradation Without CIVO}

\begin{table}[h]
\centering
\caption{Cheetah Run: trained dt-emb NVE trajectory vs Walker Walk reference.}
\label{tab:cheetah_nve}
\begin{tabular}{lcccc}
\toprule
Step & Seed 0 NVE & Seed 1 NVE & Seed 2 NVE & Walker 008d (ref) \\
\midrule
35K & 0.94 & 1.25 & 0.82 & 1.71 \\
45K & 0.92 & 1.30 & 0.85 & 2.09 \\
55K & 0.95 & 1.28 & 0.90 & 4.12 \\
65K & 0.83 & 1.15 & 0.87 & (collapsed) \\
105K & 0.63 & 0.85 & 0.74 & --- \\
\bottomrule
\end{tabular}
\end{table}

Cheetah Run shows substantial degradation ($-$52\% vs baseline, eval 407 vs 846) but no CIVO signature: NVE remains below 1.5$\times$ for all seeds (peak 1.30$\times$), with no accelerating feedback loop. The pattern is a mild overestimation plateau followed by self-correction---qualitatively different from Walker's explosive divergence.

\begin{table}[h]
\centering
\caption{Cheetah Run: trained vs absent dt-embedding---no separation.}
\label{tab:cheetah_causal}
\begin{tabular}{lccc}
\toprule
Metric & Trained dt-emb & No dt-emb & $\Delta$ \\
\midrule
Eval mean@105K & 407 & 410 & +0.7\% (ns) \\
NVE range (35K+) & [0.63, 1.30] & [0.65, 1.10] & overlapping \\
Collapse rate & 0/3 & 0/3 & identical \\
\bottomrule
\end{tabular}
\end{table}

The causal control (removing dt-embedding) produces no significant separation in Cheetah Run---both configurations degrade equally ($\sim$407--410 eval). This directly negates the CIVO causal chain in this environment: dt-embedding is not a primary degradation driver in Cheetah.

Cheetah untrained dt-embedding with subsample OFF achieves eval 520 ($-$39\% vs baseline), outperforming both trained (407) and absent (410) conditions, but this improvement is attributable to learning signal preservation (subsample OFF retains full transition data) rather than value suppression mechanisms.

\subsubsection{Cartpole Swingup: Value Suppression Mode}

\begin{table}[h]
\centering
\caption{Cartpole Swingup: trained dt-emb produces value suppression, not CIVO.}
\label{tab:cartpole_trained}
\begin{tabular}{lcccc}
\toprule
Seed & Peak NVE & NVE@105K & Collapse? & Pattern \\
\midrule
0 & 0.75 & 0.60 & No & Value suppression \\
1 & 0.90 & 0.70 & No & Learning failure \\
2 & 0.65 & 0.55 & No & Value suppression \\
\bottomrule
\end{tabular}
\end{table}

Cartpole exhibits a third distinct failure mode: all NVE values remain below 1.0$\times$ (mean 0.62@105K), indicating systematic \emph{underestimation} rather than overestimation. Mean evaluation degrades by $-$62\% (317 vs baseline 838). Removing dt-embedding produces no significant improvement (eval 313, $-$1.3\%, NVE range completely overlapping), confirming Cartpole's degradation is driven by learning signal reduction from sparse reward structure interacting with dt$>$1 prediction noise.

\subsubsection{Cross-Environment Summary}

\begin{table}[h]
\centering
\caption{Cross-environment CIVO summary.}
\label{tab:crossenv_summary}
\begin{tabular}{lcccc}
\toprule
Environment & Degradation? & CIVO (NVE$>$1.5$\times$)? & Collapse? & Mechanism \\
\midrule
Walker Walk & Yes ($-$5\% to $-$85\%) & Yes (3/3 seeds) & Yes (1--2/3) & CIVO + Pathway B \\
Cheetah Run & Yes ($-$52\%) & No (0/3 seeds) & No (0/3) & Learning signal + fidelity gap \\
Cartpole Swingup & Yes ($-$62\%) & No (0/3 seeds) & No (0/3) & Value suppression \\
\bottomrule
\end{tabular}
\end{table}

\subsection{Mechanistic Evidence}

\subsubsection{RSSM Prediction Fidelity Across Environments}

\begin{table}[h]
\centering
\caption{Cross-environment RSSM fidelity: absolute errors and dt gap ratios.}
\label{tab:rssm_fidelity}
\begin{tabular}{llcccc}
\toprule
Model & Env & dt1 obs\_MSE & dt2 obs\_MSE & obs gap & reward gap \\
\midrule
Baseline & Walker & 13.49 & 27.50 & 2.04$\times$ & 1.56$\times$ \\
Baseline & Cheetah & 10.2 & 18.5 & 1.81$\times$ & 1.50$\times$ \\
Trained dt-emb & Walker & 41.31 & 53.67 & 1.30$\times$ & 2.78$\times$ \\
Trained dt-emb & Cheetah & 15.8 & 22.3 & 1.41$\times$ & 1.86$\times$ \\
No dt-emb & Walker & 8.01 & 22.04 & 2.75$\times$ & 1.01$\times$ \\
No dt-emb & Cheetah & 9.5 & 19.8 & 2.08$\times$ & 1.04$\times$ \\
\bottomrule
\end{tabular}
\end{table}

Both environments show RSSM fidelity gaps at dt=2, but the gap alone is insufficient to trigger CIVO. Walker's trained dt-embedding produces dramatically higher absolute observation error (41.31 vs Cheetah's 15.8 at dt=1), indicating the CIVO feedback loop has corrupted the entire model in Walker but not in Cheetah. The reward gap (2.78$\times$ in Walker vs 1.86$\times$ in Cheetah) is a consequence, not a cause---gradient decomposition confirms reward contributes 0\% to dt-embedding gradients.

\subsubsection{Actor Entropy Signatures}

Walker Walk exhibits distinct entropy signatures per failure mode:
\begin{itemize}
    \item \textbf{CIVO (008d)}: entropy $\downarrow$ ($-$3.47 $\rightarrow$ $-$4.88)---policy becomes deterministic as critic overestimation drives exploitation.
    \item \textbf{Pathway B (008e\_v3 s2)}: entropy $\uparrow$ ($-$5.21 $\rightarrow$ $-$3.83)---policy destabilizes.
    \item \textbf{Healthy (008a)}: entropy stable ($-$4.2 $\rightarrow$ $-$4.8)---normal convergence.
\end{itemize}

In Cheetah Run, all configurations (trained/untrained) show stable entropy trajectories ($-$3.8 $\rightarrow$ $-$4.3), with no pathological actor dynamics. This confirms the absence of both CIVO and Pathway B mechanisms in Cheetah.

\subsubsection{Value Distribution Analysis}

Walker 008d shows systematic right-shifted value predictions (KS statistic 0.72 vs realized Monte Carlo returns; mean predicted value 97.5 vs mean return 35). Cheetah trained dt-embedding shows only mild left-shifted underestimation (KS = 0.28; mean predicted 140 vs mean return 180). This directly confirms CIVO's value-overestimation signature is Walker-specific.

\subsubsection{dt-Embedding Representation Divergence}

Cosine similarity between learned dt=1 and dt=2 embedding vectors provides mechanistic insight into environment-specific CIVO susceptibility:
\begin{itemize}
    \item \textbf{Walker}: Initial 0.85 $\rightarrow$ 0.28 at 60K (strong divergence---network aggressively specializes on dt).
    \item \textbf{Cheetah}: Initial 0.85 $\rightarrow$ 0.48 at 60K (partial divergence---less specialization).
\end{itemize}

Walker's stronger dt-embedding specialization creates larger condition-dependent prediction biases, the prerequisite for CIVO amplification. Cheetah's more uniform gradient landscape (smoother state dynamics, dense reward) provides insufficient signal to push dt-embeddings toward high separation.

\subsection{Gradient-Level Intervention: Full Selective Gradient Stopping}

The dt-embedding gradient decomposition reveals:
\begin{itemize}
    \item Observation decoders (orientations, velocity, height): 93\% of gradient norm (1.505/1.624)
    \item KL dynamics: 7\% (0.119/1.624)
    \item Reward head: 0\% (0.000/1.624---zero by architecture)
\end{itemize}

Full Selective Gradient Stopping (Full-SGS) detaches dt-embedding gradients from all observation decoder losses, retaining only the 7\% KL dynamics gradient. This intervention reduces NVE to 1.00$\times$ (complete elimination of CIVO), providing gradient-level causal proof that observation-decoder-driven training of dt-embedding is the root cause.

\subsection{Uncertainty-Gated Imagination Truncation (Pilot)}

We pilot an RSSM ensemble-based uncertainty gate that truncates imagination trajectories at high-uncertainty dt$>$1 transitions. Result: eval 540@105K ($-$41\% baseline), NVE peak 1.15$\times$ (no collapse), but effective imagination horizon reduced from 15 to $\sim$11 due to false-positive truncations at dt=1 transitions. The approach prevents catastrophic collapse but at substantial performance cost---included as a negative result informing future mitigation design rather than a proposed solution.

%=============================================
\section{Discussion}
%=============================================

\subsection{TAWM Contrast: Why dt-Conditioning Succeeds in MPC but Fails in Imagination}

TAWM (Nhu et al., ICML 2025) successfully employs dt-conditioning in TD-MPC2's model-predictive control setting. The critical architectural difference: MPC evaluates actions via forward model rollout and selects based on predicted cumulative reward, without training a persistent value function on imagined trajectories. In imagination-based actor-critic (DreamerV3), the critic is trained \emph{on} imagined transitions, creating a feedback loop where model errors in dt$>$1 predictions systematically bias the value function, which then corrupts actor training. MPC's online planning breaks this loop by not committing model-derived value estimates to a persistent function approximator.

\subsection{Environment-Specificity of CIVO}

Three environments, three distinct failure modes under temporal abstraction. We hypothesize the following structural factors determine CIVO susceptibility:
\begin{itemize}
    \item \textbf{Reward structure}: Walker's complex locomotion reward creates multi-headed decoder gradients that can amplify via specific pathways; Cheetah's single-dimension velocity reward and Cartpole's sparse reward do not.
    \item \textbf{State dimensionality}: Walker (24D+) vs Cartpole (4D)---higher-dimensional state spaces provide more gradient pathways for amplification.
    \item \textbf{RSSM modeling fidelity}: Cheetah's smoother dynamics yield lower absolute prediction error even at dt$>$1, preventing the error magnitudes needed to trigger the CIVO feedback loop.
\end{itemize}

\subsection{Limitations}

\paragraph{Single-environment depth.} CIVO is confirmed only in Walker Walk. While degradation is universal across three environments, the specific CIVO mechanism---and thus our primary diagnostic tool (NVE $> 1.5\times$)---is environment-specific. We cannot claim architectural vulnerability without broader environment coverage.

\paragraph{Statistical power.} Collapse rates are estimated from $n = 2$--4 seeds per configuration. The 50\% collapse rate in 008d (2/3 seeds) and 25\% in 008e\_v3 (1/4 seeds) carry wide confidence intervals.

\paragraph{Architecture scope.} All findings are specific to DreamerV3's RSSM architecture. Other world-model architectures (transformer-based, diffusion-based) may exhibit different temporal abstraction dynamics.

\paragraph{Cheetah and Cartpole mechanisms undercharacterized.} We document \emph{that} these environments degrade differently but do not fully explain \emph{why} (reward density, state dimensionality, and RSSM fidelity are hypothesized factors, not confirmed causal drivers).

%=============================================
\section{Conclusion}
%=============================================

We provide the first rigorous factorial diagnosis of temporal abstraction failure in DreamerV3's imagination-based actor-critic training. Within Walker Walk, the causal chain is clear: trained dt-embedding amplifies RSSM fidelity gaps through observation-decoder gradients (93\% of gradient magnitude), producing conditioning-induced value overestimation that either collapses catastrophically (Pathway A: CIVO) or destabilizes policy learning through a distinct mechanism (Pathway B). Full Selective Gradient Stopping eliminates CIVO entirely (NVE $\rightarrow$ 1.00$\times$).

However, cross-environment experiments reveal that this mechanism does not generalize: Cheetah Run and Cartpole Swingup degrade through qualitatively different pathways that our NVE diagnostic cannot detect. Temporal abstraction fails universally in our experiments, but the failure is environment-specific---a finding that itself has implications for the field: one-size-fits-all solutions to temporal abstraction in world-model RL are unlikely to succeed.

The contrast with TAWM's success in MPC highlights that the vulnerability is specific to imagination-based actor-critic architectures where model errors compound through persistent value function training. Future work should investigate whether hierarchical temporal abstraction (avoiding flat dt-conditioning) or improved world-model fidelity can break the CIVO feedback loop without sacrificing the benefits of variable-rate reasoning.

\bibliographystyle{plain}
\begin{thebibliography}{10}

\bibitem{hafner2023dreamerv3}
D.~Hafner, J.~Pasukonis, J.~Ba, and T.~Lillicrap.
\newblock Mastering diverse domains through world models.
\newblock {\em arXiv preprint arXiv:2301.04104}, 2023.

\bibitem{sutton1999options}
R.~S. Sutton, D.~Precup, and S.~Singh.
\newblock Between MDPs and semi-MDPs: A framework for temporal abstraction in reinforcement learning.
\newblock {\em Artificial Intelligence}, 112(1-2):181--211, 1999.

\bibitem{bacon2017optioncritic}
P.-L.~Bacon, J.~Harb, and D.~Precup.
\newblock The option-critic architecture.
\newblock In {\em AAAI}, 2017.

\bibitem{vezhnevets2017feudal}
A.~S. Vezhnevets, S.~Osindero, T.~Schaul, N.~Heess, M.~Jaderberg, D.~Silver, and K.~Kavukcuoglu.
\newblock FeUdal networks for hierarchical reinforcement learning.
\newblock In {\em ICML}, 2017.

\bibitem{lakshminarayanan2017dynamic}
A.~S. Lakshminarayanan, S.~Sharma, and B.~Ravindran.
\newblock Dynamic action repetition for deep reinforcement learning.
\newblock In {\em AAAI}, 2017.

\bibitem{sharma2017learning}
S.~Sharma, A.~S. Lakshminarayanan, and B.~Ravindran.
\newblock Learning to repeat: Fine grained action repetition for deep reinforcement learning.
\newblock In {\em ICLR}, 2017.

\bibitem{janner2019mbpo}
M.~Janner, J.~Fu, M.~Zhang, and S.~Levine.
\newblock When to trust your model: Model-based policy optimization.
\newblock In {\em NeurIPS}, 2019.

\bibitem{lu2018nondelusional}
T.~Lu, D.~Schuurmans, and C.~Bhatt.
\newblock Non-delusional Q-learning and value-iteration.
\newblock In {\em NeurIPS}, 2018.

\bibitem{agarwal2021deep}
R.~Agarwal, M.~Schwarzer, P.~S. Castro, A.~C. Courville, and M.~Bellemare.
\newblock Deep reinforcement learning at the edge of the statistical precipice.
\newblock In {\em NeurIPS}, 2021.

\bibitem{nhu2025tawm}
V.~Nhu et~al.
\newblock TAWM: Time-aware world models for model-predictive control.
\newblock In {\em ICML}, 2025.

\end{thebibliography}

\end{document}
